\documentclass[12pt, twocolumn, landscape, a4paper]{article}

\usepackage[margin=1cm]{geometry}
\usepackage[utf8]{inputenc}
\usepackage[spanish]{babel}
\usepackage{amssymb, amsmath, amsbsy}
\usepackage{tasks}
\usepackage{enumerate}
\usepackage{enumitem}
\usepackage{graphicx}

\usepackage[pdftex]{hyperref}
\hypersetup{pdftitle={Fracciones},pdfauthor={Luis Carrillo Gutiérrez}}

\fontfamily{cmss}\selectfont
\renewcommand{\familydefault}{\sfdefault}
\pagestyle{empty}

\begin{document}
\subsection*{Operaciones con fracciones / Racionales ($\mathbb{Q}$)}

\paragraph*{Adición/Sustracción} --  mismo denominador
$$
\dfrac{a}{b} \pm \dfrac{c}{b} = \dfrac{a \pm c}{b}
$$

\paragraph*{Adición} -- diferentes denominadores
$$
\dfrac{a}{b} + \dfrac{c}{d} = \dfrac{(a \cdot d) + (b \cdot c)}{b \cdot d}
$$

\paragraph*{Sustracción} -- diferentes denominadores
$$
\dfrac{a}{b} - \dfrac{c}{d} = \dfrac{(a \cdot d) - (b \cdot c)}{b \cdot d}
$$

\paragraph*{Multiplicación}
$$
\Big(\dfrac{a}{b}\Big) \times \Big(\dfrac{c}{d}\Big) = \dfrac{a \cdot c}{b \cdot d}
$$

\paragraph*{División}
$$
\Big(\dfrac{a}{b}\Big) \div \Big(\dfrac{c}{d}\Big) = \dfrac{\dfrac{a}{b}}{\dfrac{c}{d}} = \dfrac{a \cdot d}{b \cdot c}
$$
\end{document}